\section{负载形态：一个前置的输入}
\label{sec:load}

前四节换了四次切法：成本链、实现平面、SLO 方向。本节要说的东西
性质不同——\textbf{负载形态不是决策维度，是决策的输入}。

理由很简单：同一套技巧、同一个 SLO，如果输入长什么样不同，
结论会翻转。而翻转的幅度大到\textbf{可以让前面整节的分析失去意义}。

\subsection{两个数字足以翻转结论}

先看一组真实 agentic 负载的统计。Codex / SWE-bench Pro 的
$610$ 条 trace，中位数每条 $33$ 轮\texttt{[c-load]}：

\begin{itemize}[leftmargin=1.3em,itemsep=1.5pt,topsep=2pt]
  \item \textbf{缓存命中率 $94.2\%$}
  \item \textbf{输入输出比 $131:1$}
  \item 每轮上下文增长约 $2{,}242$ token，中位轨迹从 $12$K 涨到 $80$K
  \item 轮间隔中位数 $5.2$\,s，P99 为 $81.4$\,s
\end{itemize}

\noindent 这两个加粗的数字各自推翻一件事。

\noindent\textbf{$94.2\%$ 命中率推翻「省传输量」的意义。} 命中意味着
KV 已经在本地，根本不需要搬——不是搬得更快，是不搬。
当命中率到九成以上，第 \ref{sec:technique} 节 A 类（逐层传输、
KV 压缩、量化）与 B 类（TBO/SBO/DBO）的收益空间被压缩到
剩下的那几个百分点里，而它们的复杂度与前置条件一分不减。

\noindent 这不是说 A、B 两类没用，而是说\textbf{它们的价值取决于
一个本文前面没有当作变量的量}。同样的技巧，在命中率 $0$ 的负载上是
主线，在命中率 $94.2\%$ 的负载上是边角。

\noindent\textbf{$131:1$ 的输入输出比推翻 P/D 配比的常规直觉。}
分离的立论基础是两阶段资源画像不同，故可以各自配最优资源。
但当输入是输出的 $131$ 倍时，\textbf{prefill 侧承担的计算量
与 decode 侧完全不成比例}，P:D 的实例比会极端倾斜。
这时「均衡负载」的含义也变了：不是两类实例之间调平，
而是少量 D 实例要服务大量 P 实例产出的 KV。

\noindent 本文第 \ref{sec:boundary} 节给出的窗口判据里有一个 I/O 比维度
（$0.01$ 到 $1{,}000$ 的扫描），那里的结论是窗口只跨一个数量级。
现在可以补一句：\textbf{真实 agentic 负载落在 $131$，正在窗口内，
但它的命中率把 A/B 两类挤到了边缘}——两个维度的作用不是相加的。

\subsection{负载形态还决定「传输是否发生」}

命中率之外，负载形态还从另一个方向影响结论：它决定
\textbf{那些被设计的传输到底会不会被触发}。

\noindent\textbf{容量与驱逐。} 一个 $100$K token 的上下文可占数 GB
存储——例如 Kimi-2.5 的 FP8 KV cache 约 $3.8$\,GB\texttt{[c-load]}。
一个繁忙实例服务多个长会话时，这些大的前缀缓存会迅速撑满本地容量
并触发驱逐\texttt{[c-load]}。此时「KV 留在本地」这个前提不成立。

\noindent\textbf{跨实例未命中。} 更微妙的是这一条：为了均衡负载，
路由器可能不会把会话的下一轮调度到同一个实例上；一旦会话被迁到
别的实例，那个实例\textbf{从未见过该前缀，必须从头重算}
\texttt{[c-load]}。

\noindent 这条与上一节的反例是同一件事的两面。上一节说
「按队列深度均衡劣于 session 粘性路由」；这里给出了它的
负载形态条件——\textbf{轮间隔中位数只有 $5.2$\,s}。间隔短，
前缀大概率还在原实例上，粘性就赢；若间隔长到缓存已被驱逐，
粘性的优势也就消失了。

\noindent 所以「该不该做粘性路由」这个问题没有独立答案，
它的答案是一个关于轮间隔与缓存容量的函数。

\subsection{命中率本身可能是假象}

这一节必须提一个测量陷阱，因为它会让人把假的负载形态当成真的。

\noindent vLLM 的 KV 块按内容的哈希寻址。而哈希链默认用
\texttt{os.urandom()} 按进程播种——两个引擎对完全相同的 token
会算出\textbf{不同的哈希}，于是每一次跨实例查找都\textbf{静默地未命中}，
\textbf{呈现出一个看起来完全合理的零命中率}\texttt{[c-load]}。

\noindent 修法是把 \texttt{PYTHONHASHSEED} 在所有共享同一区域的副本上
设为同一个值\texttt{[c-load]}。

\noindent 这条值得单独写出来的原因：它的失败模式是
「\textbf{数字合理、无报错、无告警}」。一个零命中率在诊断时会被
归因为「负载的前缀复用不好」，而不是「哈希没对齐」。
\textbf{这与本文反复强调的「不可校验 $\neq$ 通过」是同一类错误}——
一个看起来像测量的零，与一个真的零，在报表上没有区别。

\noindent 还有一层更细的粒度问题：offload 的粒度是 $256$ token
（该研究的配置），\textbf{任何不是 $256$ 整数倍的前缀都要重算尾部的
部分块}\texttt{[c-load]}。这意味着即使命中率很高，
实际省下的计算也取决于输入长度与 $256$ 的余数——
一个纯属实现细节的量，却直接改变收益。

\subsection{缓存能换回多少：三个不同性质的数}

有了负载形态，才能谈缓存的价值。表 \ref{tab:load} 汇总了
三个来源的数——它们差别很大，
而且\textbf{差别来自对比基线不同}，把它们并列会误读。

\begin{table}[t]
\centering
\caption{缓存能换回多少：三个不同性质的数。幅度差别主要来自
基线不同，故第三列是这张表的关键。}
\label{tab:load}
\small
\begin{tabular}{@{}p{2.6cm}p{1.9cm}p{4.4cm}@{}}
\toprule
来源 & 幅度 & 基线是什么 \\
\midrule
CXL 可行性研究 & TTFT 降 $5.5\times$ 至 $36.6\times$
  & 同节点复用（外部命中率 $95.4$--$99.5\%$） \\
Mooncake 集成 & 吞吐 $3.8\times$、TTFT 降 $46\times$、
  端到端 $8.6\times$ & 本地 KV offload，非「无缓存」 \\
\bottomrule
\end{tabular}
\end{table}

\noindent CXL 那项研究还报了一个有意思的派生量：\textbf{sharing gap}
（跨节点复用与同节点复用的延迟比）只有 $1\%$ 至 $4\%$\texttt{[c-load]}，
说明在它的测试床上跨节点复用几乎没有额外代价。

\noindent 但这项研究的作者自己把边界写清了，值得原样引用：两个副本
\textbf{都跑完整引擎}，因此该研究展示的是\textbf{内存分离而非
P/D 分离}，并声明它是\textbf{可行性研究而非性能评估}
\texttt{[c-load]}。

\noindent 这句话对本文的用法是个提醒：\textbf{一个数和它成立的范围是
一起交付的}。若只引「$46\times$ 更低 TTFT」，读者会以为那是分离的收益；
实际上它测的是「有个共享 KV 池 vs 只有本地 offload」。
两者不是一件事。

\subsection{这一步在判据里的位置}

写到这里，五次切法的关系可以收拢成一条使用顺序：

\begin{enumerate}[leftmargin=1.4em,itemsep=2.5pt,topsep=3pt]
  \item \textbf{先看负载形态}：命中率、I/O 比、轮间隔、上下文长度分布。
        这是输入，不是可选项——它决定后面哪些类有意义。
  \item \textbf{再定 SLO 形状}：哪个更紧、比值多少、达成率目标。
  \item \textbf{按实现平面查漏}：数据面 / 执行面 / 控制面。
  \item \textbf{最后按成本链取舍}：四类里每类选一个。
\end{enumerate}

\noindent 负载形态排在最前，理由是它可以\textbf{让后面三步的结论失效}，
而不是让它们更精确。SLO 只能翻转方向（上一节），
负载形态能直接取消一整类技巧的用武之地（本节）。

\noindent 反过来也有一个正面结论：\textbf{当命中率足够高时，
最有效的优化不是把 KV 搬得更快，而是让 KV 不必被搬}。
这又把话题带回了本文反复出现的那条主线——\textbf{收益来自
「让两阶段不必做同一件事」，而最高效的做法是让某件事根本不发生}。
分离是这条主线里最贵的一个实例；前缀命中是同一个道理在缓存层的
体现，且它比分离便宜得多。
